\documentclass[11pt,a4paper]{article}

\usepackage{amsmath,amssymb,amsthm}
\usepackage{mathtools}
\usepackage{geometry}
\usepackage{hyperref}
\usepackage{booktabs}
\usepackage{xcolor}
\usepackage{listings}
\usepackage{enumitem}
\usepackage{fancyhdr}
\usepackage{tocloft}
\usepackage{microtype}
\usepackage{parskip}

\geometry{margin=2.5cm, headheight=14pt}

\hypersetup{
  colorlinks=true,
  linkcolor=blue!60!black,
  urlcolor=blue!60!black,
  pdfauthor={Aegis Quant},
  pdftitle={Aegis Quant Pricing Engine -- Theoretical Background}
}

\lstset{
  basicstyle=\ttfamily\small,
  keywordstyle=\color{blue!70!black}\bfseries,
  commentstyle=\color{gray},
  breaklines=true,
  frame=single,
  backgroundcolor=\color{gray!8},
  language=[Sharp]C,
  showstringspaces=false
}

\newcommand{\dd}{\mathrm{d}}
\newcommand{\E}{\mathbb{E}}
\newcommand{\Q}{\mathbb{Q}}
\newcommand{\QT}{\mathbb{Q}^T}
\newcommand{\Qd}{\mathbb{Q}^d}
\newcommand{\R}{\mathbb{R}}
\newcommand{\N}{\mathcal{N}}
\newcommand{\F}{\mathcal{F}}
\newcommand{\coderef}[1]{\texttt{\color{blue!60!black}#1}}

\theoremstyle{definition}
\newtheorem{defn}{Definition}[section]
\newtheorem{prop}{Proposition}[section]
\newtheorem{remark}{Remark}[section]

\pagestyle{fancy}
\fancyhf{}
\rhead{\textit{Aegis Quant -- Theory}}
\lhead{\leftmark}
\cfoot{\thepage}

\title{\textbf{Aegis Quant Pricing Engine}\\[0.4em]
       \large Theoretical Background and Mathematical Derivations}
\author{Aegis Quant Library}
\date{June 2026}

\begin{document}
\maketitle
\tableofcontents
\newpage

%%==========================================================================
\section{Introduction}
%%==========================================================================

This document provides the mathematical foundations underlying every model,
analytical formula, and Monte Carlo simulation implemented in the Aegis
quant pricing engine.  Each section states the stochastic process, the
probability measure, the key assumptions, and the derivation of the pricing
formula, with explicit references to the C\# source files where the
corresponding implementation lives.

The library prices:
\begin{itemize}[nosep]
  \item Vanilla and exotic FX options (Garman-Kohlhagen, SABR, barriers, Asians, digitals, baskets)
  \item Vanilla and exotic equity options (Black-Scholes-Merton, barriers, Asians, digitals)
  \item SABR-based barrier options under the T-forward measure
\end{itemize}

%%==========================================================================
\section{Probability Measures and the Change-of-Numeraire Theorem}
\label{sec:measures}
%%==========================================================================

\subsection{Numeraires and equivalent martingale measures}

Let $(\Omega,\F,\{\F_t\},P)$ be a filtered probability space satisfying the
usual conditions.  A \emph{numeraire} is any strictly positive $\F_t$-adapted
price process $N(t)$.

\begin{defn}[Equivalent martingale measure (EMM)]
A probability measure $\Q\sim P$ is an EMM associated with numeraire $N$ if
every asset price $V(t)$ (after dividing by $N$) is a local $\Q$-martingale:
\[
  \frac{V(t)}{N(t)} = \E^\Q\!\left[\frac{V(T)}{N(T)}\,\Big|\,\F_t\right].
\]
\end{defn}

Absence of arbitrage is equivalent (under integrability conditions) to the
existence of at least one EMM.

\subsection{The risk-neutral measure $\Q$}

The most common choice of numeraire is the \emph{domestic money-market account}
$B(t)=\exp\!\bigl(\int_0^t r(s)\,\dd s\bigr)$ with $r$ the short rate.  The
associated EMM is the \emph{risk-neutral measure} $\Q$.  Under a flat rate $r$:
\[
  V(0) = e^{-rT}\,\E^\Q[V(T)].
\]

For FX, the domestic money-market account $B^d$ denominates prices in the
domestic currency (e.g.\ USD); for equity it is the risk-free bank account.
Code reference: \coderef{GbmOptionMCPricer.cs} (\texttt{DiscountFactor}
computed as $e^{-r_d T}$).

\subsection{The T-forward measure $\QT$}

The \emph{T-forward measure} (or \emph{T-annuity measure}) uses the zero-coupon
domestic bond $P(t,T)$ as numeraire.  By the change-of-numeraire theorem, if
$\Q$ is the risk-neutral measure:
\[
  \frac{\dd\QT}{\dd\Q}\bigg|_{\F_T}
  = \frac{P(0,T)^{-1}}{P(T,T)^{-1}}\cdot\frac{B(0)}{B(T)}
  = \frac{P(0,T)^{-1}}{1}\cdot e^{-rT}
  = 1 \quad\text{(flat rates, }P(0,T)=e^{-rT}\text{)}.
\]

Under flat rates the two measures coincide \emph{for terminal payoffs}, so
\[
  V(0) = P(0,T)\,\E^T[V(T)].
\]

The key property is that any traded forward $F(t,T)=S(t)/P(t,T)$ is a
$\QT$-martingale.  This makes the $\QT$-measure the natural setting for
pricing futures and forward-settled payoffs: no drift term appears in the
forward's SDE.  Code reference: \coderef{FXSABRMCPricer.cs} (SABR simulated
under $\QT$; \texttt{DiscountFactor}$=e^{-r_d T}$,
\texttt{InitialForward}$=S_0 e^{(r_d-r_f)T}$).

%%==========================================================================
\section{Geometric Brownian Motion and the It\^{o} Formula}
\label{sec:gbm}
%%==========================================================================

\subsection{The GBM SDE and its solution}

The spot process $S$ satisfies under $\Q$:
\[
  \frac{\dd S}{S} = \mu\,\dd t + \sigma\,\dd W^\Q,
\]
where $W^\Q$ is a $\Q$-Brownian motion.  Applying It\^{o}'s lemma to
$X=\ln S$:
\[
  \dd X = \dd\ln S = \frac{1}{S}\,\dd S - \frac{1}{2S^2}\,(\dd S)^2
  = \Bigl(\mu - \tfrac{1}{2}\sigma^2\Bigr)\dd t + \sigma\,\dd W^\Q.
\]
Integrating:
\[
  \boxed{S(T) = S(0)\exp\!\Bigl[\Bigl(\mu-\tfrac{1}{2}\sigma^2\Bigr)T + \sigma W^\Q(T)\Bigr].}
\]
The $-\frac{1}{2}\sigma^2$ correction is the \emph{It\^{o} correction}
(convexity adjustment).  Under $\Q$ one sets $\mu = r_d - r_f$ (FX) or
$\mu = r - q$ (equity).

\subsection{Log-Euler discretization}

Applying the exact formula over a time step $[t_i,t_{i+1}]$ of length
$\Delta t$:
\[
  S(t_{i+1}) = S(t_i)\exp\!\Bigl[\bigl(f_d(t_i,t_{i+1})-f_c(t_i,t_{i+1})-\tfrac{1}{2}\sigma^2\bigr)\Delta t
  + \sigma\sqrt{\Delta t}\,Z_i\Bigr],
  \quad Z_i\sim\N(0,1),
\]
where $f_d,f_c$ are the instantaneous forward rates from the respective curves.
This is \emph{exact} for GBM (no discretization error), which is why the
Aegis MC pricers use log-Euler rather than plain Euler.

Code reference: \coderef{GbmOptionMCPricer.cs}, method
\texttt{SimulateSpotPath}.

%%==========================================================================
\section{The FX Forward Measure and the Driftless Forward}
\label{sec:fx_forward}
%%==========================================================================

\subsection{FX covered-interest parity}

Let $S(t)$ be the spot FX rate (domestic per unit foreign, e.g.\ USD/EUR).
The foreign money-market account $B^f(t)=e^{r_f t}$ is a domestic asset with
price $S(t)B^f(t)$ in domestic currency.  Absence of arbitrage requires:
\[
  F(0,T) = S(0)\,\frac{B^d(T)}{B^f(T)} = S(0)\,e^{(r_d-r_f)T}.
\]
This is \emph{covered-interest parity}.  Code:
\coderef{FXSABRMCPricer.cs} (\texttt{InitialForward}), and
\coderef{GbmOptionMCPricer.cs} (implicit via drift $r_d - r_f$).

\subsection{Derivation: the forward is driftless under $\QT$}

Under the domestic risk-neutral measure $\Qd$:
\[
  \frac{\dd S}{S} = (r_d - r_f)\,\dd t + \sigma\,\dd W^d.
\]
The T-forward is $F(t,T)=S(t)/P^d(t,T)$.  Under flat rates
$P^d(t,T)=e^{-r_d(T-t)}$, so $F(t,T)=S(t)e^{r_d(T-t)}\cdot e^{-r_d(T-t)+r_f(T-t)}\cdots$
More precisely, applying It\^{o}'s lemma to $F=S\,e^{(r_d-r_f)(T-t)}$ under
$\Qd$:

\begin{align*}
  \dd F &= e^{(r_d-r_f)(T-t)}\dd S - (r_d-r_f)\,S\,e^{(r_d-r_f)(T-t)}\dd t \\
        &= e^{(r_d-r_f)(T-t)}\bigl[(r_d-r_f)S\,\dd t + \sigma S\,\dd W^d\bigr]
           -(r_d-r_f)S\,e^{(r_d-r_f)(T-t)}\dd t \\
        &= \sigma F\,\dd W^d.
\end{align*}

The drift cancels exactly.  This confirms $F(t,T)$ is a $\Qd$-local
martingale (and hence a $\QT$-martingale by the equivalence of the two
measures under flat rates).  Under $\QT$:
\[
  \boxed{\frac{\dd F}{F} = \sigma\,\dd W^T,\qquad \text{(no drift term).}}
\]

This is why \coderef{FXSABRMCPricer.cs} sets the drift to zero for the
forward process: ``\emph{Under the T-forward measure, the forward F is a
martingale --- no drift term.}''

\subsection{Pricing FX options in each measure}

The same European call price can be written in two equivalent ways:
\begin{align*}
  C &= e^{-r_d T}\,\E^d[\max(S(T)-K,0)] \tag{risk-neutral on $S$}\\
    &= e^{-r_d T}\,\E^T[\max(F(T)-K,0)] \tag{T-forward on $F$, same terminal value}
\end{align*}
Both representations are used in Aegis:
\begin{itemize}[nosep]
  \item \coderef{GarmanKohlhagen.cs} uses the risk-neutral representation (on $S$).
  \item \coderef{FXSABRMCPricer.cs} and \coderef{FXSABRVanillaOptionMCPricer.cs} use the T-forward representation (on $F$).
  \item \coderef{BlackForward.cs} prices the T-forward payoff without discounting.
\end{itemize}

%%==========================================================================
\section{Garman-Kohlhagen Formula}
\label{sec:gk}
%%==========================================================================

\subsection{Model}

Under the domestic risk-neutral measure $\Qd$:
\[
  \frac{\dd S}{S} = (r_d - r_f)\,\dd t + \sigma\,\dd W^d.
\]

\subsection{Formula derivation}

By the risk-neutral pricing formula:
\[
  C = e^{-r_d T}\,\E^d[\max(S(T)-K,0)].
\]
Write $S(T)=S_0\exp\bigl[(r_d-r_f-\tfrac{1}{2}\sigma^2)T+\sigma\sqrt{T}\,Z\bigr]$ with $Z\sim\N(0,1)$.
Let $d_1,d_2$ be defined by the boundary $S(T)=K$ in $Z$-space:
\[
  d_1 = \frac{\ln(S_0/K)+(r_d-r_f+\tfrac{1}{2}\sigma^2)T}{\sigma\sqrt{T}},
  \qquad d_2 = d_1 - \sigma\sqrt{T}.
\]
A standard Gaussian integration (splitting the expectation into
``asset term'' and ``cash term'') gives:
\[
  \boxed{C = S_0\,e^{-r_f T}\,\Phi(d_1) - K\,e^{-r_d T}\,\Phi(d_2),}
\]
\[
  P = K\,e^{-r_d T}\,\Phi(-d_2) - S_0\,e^{-r_f T}\,\Phi(-d_1).
\]

Code: \coderef{GarmanKohlhagen.cs}, \texttt{Price()}.

\textbf{Edge case.} At $\sigma=0$ the path is deterministic, $S(T)=F=S_0e^{(r_d-r_f)T}$,
and the price collapses to the discounted intrinsic:
$e^{-r_d T}\max(\phi(F-K),0)$.

%%==========================================================================
\section{Black-Scholes-Merton Formula (Equity)}
\label{sec:bsm}
%%==========================================================================

The equity model is structurally identical to Garman-Kohlhagen with the
substitutions $r_f\to q$ (dividend yield) and $r_d\to r$:
\[
  \frac{\dd S}{S} = (r - q)\,\dd t + \sigma\,\dd W,
\]
giving:
\[
  C = S_0\,e^{-qT}\,\Phi(d_1) - K\,e^{-rT}\,\Phi(d_2),
  \quad d_1=\frac{\ln(S_0/K)+(r-q+\tfrac{1}{2}\sigma^2)T}{\sigma\sqrt{T}}.
\]

Code: \coderef{BlackScholes.cs}, \texttt{Price()};
\coderef{GbmOptionMCPricer.cs} (Monte Carlo path generation with
$r_{\mathrm{discount}}=r$, $r_{\mathrm{carry}}=q$).

%%==========================================================================
\section{Black Forward Formula}
\label{sec:black_fwd}
%%==========================================================================

Under the T-forward measure $\QT$ the forward satisfies $\dd F/F=\sigma\,\dd W^T$
(driftless), with $F(T)=F_0\exp\bigl(-\tfrac{1}{2}\sigma^2 T+\sigma\sqrt{T}\,Z\bigr)$.
The undiscounted European payoff:
\[
  C_{\rm fwd} = \E^T[\max(F(T)-K,0)]
              = F_0\,\Phi(d_1^F) - K\,\Phi(d_2^F),
\]
\[
  d_1^F = \frac{\ln(F_0/K)+\tfrac{1}{2}\sigma^2 T}{\sigma\sqrt{T}},
  \qquad d_2^F = d_1^F - \sigma\sqrt{T}.
\]
The full price is $P(0,T)\cdot C_{\rm fwd}$.  Note $d_1^F$ differs from
GK's $d_1$ only by the forward substitution $F_0\equiv S_0 e^{(r_d-r_f)T}$ and
the absence of the rate terms inside the logarithm.

Code: \coderef{BlackForward.cs}, \texttt{Call()}, \texttt{Put()}.
Used as the lognormal-SABR ($\beta=1$, $\nu>0$) benchmark in
\coderef{FXSABRVanillaOptionMCPricerTests.cs}.

%%==========================================================================
\section{Bachelier (Normal) Model}
\label{sec:bachelier}
%%==========================================================================

The normal model replaces the lognormal GBM with arithmetic Brownian motion:
\[
  \dd F = \sigma_N\,\dd W^T,
\]
so $F(T) = F_0 + \sigma_N\sqrt{T}\,Z$.  The European option price is:
\[
  \boxed{C = P(0,T)\bigl[(F_0-K)\,\Phi(d) + \sigma_N\sqrt{T}\,\phi(d)\bigr],}
\]
\[
  P_{\rm bch} = P(0,T)\bigl[(K-F_0)\,\Phi(-d) + \sigma_N\sqrt{T}\,\phi(d)\bigr],
\]
where $d=(F_0-K)/(\sigma_N\sqrt{T})$ and $\phi$ is the standard normal PDF.

This can be derived by direct Gaussian integration:
\begin{align*}
  C &= P(0,T)\int_{(K-F_0)/(\sigma_N\sqrt{T})}^{\infty}
       \bigl(F_0+\sigma_N\sqrt{T}\,z-K\bigr)\phi(z)\,\dd z.
\end{align*}
Splitting: $(F_0-K)\Phi(-d_{\rm lower})+\sigma_N\sqrt{T}\,\phi(d_{\rm lower})$
where $d_{\rm lower}=-d$, and using $\phi(d)=\phi(-d)$ gives the stated formula.

\textbf{Usage in Aegis.} $\sigma_N$ is the \emph{normal volatility} (in
forward units per $\sqrt{\text{year}}$).  Code: \coderef{Bachelier.cs},
\texttt{Price()}.  Used as the analytical benchmark for normal-SABR ($\beta=0$)
in \coderef{FXSABRVanillaOptionMCPricerTests.cs}.

%%==========================================================================
\section{Digital Options}
\label{sec:digital}
%%==========================================================================

\subsection{Cash-or-nothing}

Pays a fixed amount $Q$ if $S(T)>K$ (call) or $S(T)<K$ (put):
\[
  \mathrm{CoN\_Call} = e^{-rT}\,Q\,\Phi(d_2),
  \qquad
  \mathrm{CoN\_Put}  = e^{-rT}\,Q\,\Phi(-d_2).
\]
This follows immediately from $\E^Q[\mathbf{1}_{S(T)>K}]=\Phi(d_2)$, which
is the ``cash term'' of Black-Scholes.

\subsection{Asset-or-nothing}

Pays $S(T)$ if $S(T)>K$:
\[
  \mathrm{AoN\_Call} = S_0\,e^{-qT}\,\Phi(d_1).
\]
This is the ``asset term'' of Black-Scholes; equivalently, it equals the
derivative of the call with respect to $K$ (up to sign), giving the
connection $\partial C/\partial K = -e^{-rT}\Phi(d_2)$ (call spread / cash
digital).

Code: \coderef{DigitalBlackScholes.cs}, \texttt{CashOrNothing()},
\texttt{AssetOrNothing()}.

%%==========================================================================
\section{Geometric Asian Options}
\label{sec:asian}
%%==========================================================================

\subsection{Model}

Monitoring at $n$ equally-spaced dates $t_i=iT/n$, $i=1,\ldots,n$.
Under $\Q$, $\ln S(t_i)$ is Gaussian:
\[
  \ln S(t_i) = \ln S(0) + \bigl(r-q-\tfrac{1}{2}\sigma^2\bigr)t_i
               + \sigma\,W(t_i).
\]

\subsection{Moments of the geometric average}

The discrete geometric average is:
\[
  G = \exp\!\Bigl(\tfrac{1}{n}\textstyle\sum_{i=1}^n \ln S(t_i)\Bigr).
\]
Since $\ln S(t_i)$ are jointly Gaussian, $\ln G$ is also Gaussian.

\textbf{Mean:}
\[
  \E[\ln G] = \ln S(0) + \bigl(r-q-\tfrac{1}{2}\sigma^2\bigr)\frac{T}{n}
              \cdot\frac{n(n+1)/2}{n}
  = \ln S(0) + \bigl(r-q-\tfrac{1}{2}\sigma^2\bigr)T\,\frac{n+1}{2n}.
\]

\textbf{Variance:}
\begin{align*}
  \mathrm{Var}[\ln G]
  &= \frac{\sigma^2}{n^2}\sum_{i=1}^n\sum_{j=1}^n \mathrm{Cov}\bigl(W(t_i),W(t_j)\bigr)\\
  &= \frac{\sigma^2}{n^2}\sum_{i,j}\min(t_i,t_j)
   = \frac{\sigma^2}{n^2}\cdot\frac{T}{n}\cdot\frac{n(n+1)(2n+1)}{6}\\
  &= \sigma^2 T\,\frac{(n+1)(2n+1)}{6n^2}.
\end{align*}
(using $\sum_{i=1}^n\sum_{j=1}^n\min(i,j)=n(n+1)(2n+1)/6$).

\subsection{Reduction to Black formula}

Since $G$ is lognormal, the Asian price is exactly the Black formula on a
``synthetic'' forward $F_G$ with effective vol $\sigma_{\rm eff}$:
\[
  F_G = S(0)\exp\!\Bigl(\E[\ln G]-\ln S(0)+\tfrac{1}{2}\mathrm{Var}[\ln G]\Bigr),
  \quad \sigma_{\rm eff}=\sqrt{\mathrm{Var}[\ln G]/T}.
\]
\[
  \mathrm{Asian\_Call} = e^{-rT}\,\mathrm{BlackFwd}(F_G,K,T,\sigma_{\rm eff}).
\]

Code: \coderef{GeometricAsian.cs}, \texttt{Price()}.  The arithmetic-average
Asian (no closed form) is computed by Monte Carlo in
\coderef{EQAsianOptionMCPricer.cs} and \coderef{FXAsianOptionMCPricer.cs}.

%%==========================================================================
\section{Zero Curves and Forward Rates}
\label{sec:curves}
%%==========================================================================

\subsection{Discount factor and zero rate}

A zero-coupon bond with maturity $T$ is priced:
\[
  P(0,T) = e^{-r(T)\cdot T},
\]
where $r(T)$ is the continuously-compounded zero rate.  The function
\texttt{ZeroCurve.InterpolateRate} performs linear interpolation in the
zero-rate space over the \texttt{Pillars} term structure.

Code: \coderef{ZeroCurve.cs}, \texttt{InterpolateRate()}, \texttt{DiscountFactor()}.

\subsection{Instantaneous forward rate}

The instantaneous (continuously-compounded) forward rate between $t_0$ and $t_1$:
\[
  f(t_0,t_1) = \frac{r(t_1)\cdot t_1 - r(t_0)\cdot t_0}{t_1-t_0}.
\]
This ensures: $e^{-r(t_0)t_0}\cdot e^{-f(t_0,t_1)(t_1-t_0)}=e^{-r(t_1)t_1}$,
i.e.\ the forward discount is consistent with the zero curve.

Code: \coderef{ZeroCurve.cs}, \texttt{ForwardRate()}.  Used in
\coderef{GbmOptionMCPricer.cs} to compute the per-step drift from the full
curve:
\[
  \mu_i\cdot\Delta t = \bigl[f_d(t_i,t_{i+1})-f_c(t_i,t_{i+1})-\tfrac{1}{2}\sigma^2\bigr]\Delta t.
\]

%%==========================================================================
\section{Discrete Dividends: the Escrowed Jump Model}
\label{sec:dividends}
%%==========================================================================

\subsection{Model}

A cash dividend $D_i$ is paid on ex-date $\tau_i\in(0,T]$.  The equity spot
drops by the dividend amount at $\tau_i$ (the \emph{jump} or \emph{escrowed}
model):
\[
  S(\tau_i^+) = S(\tau_i^-) - D_i.
\]
Combined with the continuous GBM drift:
\[
  S(t_{k+1}) = S(t_k)\exp\!\bigl[\mu_k\Delta t + \sigma\sqrt{\Delta t}\,Z_k\bigr]
               - D_k,
\]
where $D_k = \sum_{\tau_i\in(t_k,t_{k+1}]}D_i$ is the sum of all dividends
ex-dating during step $k$.

\subsection{Validation}

The model requires $\sum_{i:\tau_i\le T}D_i < S(0)$ (otherwise the forward
becomes negative in the deterministic limit).  Code:
\coderef{GbmOptionMCPricer.cs}, \texttt{PrepareDividends()}.

%%==========================================================================
\section{The SABR Stochastic Volatility Model}
\label{sec:sabr}
%%==========================================================================

\subsection{Model dynamics (T-forward measure)}

Hagan, Kumar, Lesniewski, Woodward (2002).  Under the T-forward measure $\QT$:
\begin{align}
  \dd F &= \sigma\,C(F)\,\dd W_1^T, \label{eq:sabr_F}\\
  \dd\sigma &= \nu\,\sigma\,\dd W_2^T, \label{eq:sabr_sigma}\\
  \dd W_1^T\,\dd W_2^T &= \rho\,\dd t. \label{eq:sabr_corr}
\end{align}

The CEV function $C(F)$ controls the backbone:
\begin{itemize}[nosep]
  \item $\beta=0$: $C(F)=1$ (normal/additive diffusion, forward can be negative)
  \item $\beta=1$: $C(F)=F$ (lognormal/multiplicative diffusion)
  \item $0<\beta<1$: $C(F)=F^\beta$ (CEV, absorbing boundary at $F=0$)
\end{itemize}

Initial conditions: $F(0)=F_0=S_0\exp\bigl[(r_d-r_f)T\bigr]$, $\sigma(0)=\alpha$.
The option price is:
\[
  V = P(0,T)\,\E^T[\max(\phi(F(T)-K),0)],\quad\phi=\pm 1.
\]

Code: \coderef{FXSABRMCPricer.cs} (abstract base); all parameters in
\coderef{SabrParameters} (protobuf).

\subsection{Why the forward is a martingale under $\QT$}

Since $C(F)\geq 0$ and the SDE \eqref{eq:sabr_F} has no drift term, $F$ is a
local $\QT$-martingale by construction (the T-forward measure was chosen
precisely for this purpose; see Section~\ref{sec:fx_forward}).

\subsection{Euler--Maruyama and log-Euler discretization}

\subsubsection{Volatility process (log-Euler, exact)}

Since $\sigma$ follows lognormal dynamics $\dd\sigma=\nu\sigma\,\dd W_2^T$,
the exact solution over $[t,t+\Delta t]$ is:
\[
  \sigma(t+\Delta t)=\sigma(t)\exp\!\bigl[\nu\sqrt{\Delta t}\,\tilde W_2
                                          -\tfrac{1}{2}\nu^2\Delta t\bigr],
\]
where $\tilde W_2$ is a standard normal.  Log-Euler for $\sigma$ is exact
because the SDE is linear-multiplicative in $\sigma$ with no $F$-dependence.

Code: \coderef{FXSABRMCPricer.cs}, \texttt{SimulateSabrPath()}.

\subsubsection{Brownian decomposition}

Two independent factors $Z_1,Z_2\sim\N(0,1)$ from the simulation cube are
combined to reproduce the correlation $\rho$:
\[
  \tilde W_1 = Z_1,\qquad \tilde W_2 = \rho Z_1 + \sqrt{1-\rho^2}\,Z_2.
\]
This factorization keeps $\rho$ as a pure model parameter; bumping $\rho$
never requires regenerating the cube.

Code: \coderef{FXSABRMCPricer.cs}, field \texttt{\_rhoComplement}$=\sqrt{1-\rho^2}$.

\subsubsection{Forward process}

\textbf{$\beta=1$ (lognormal, log-Euler, exact for frozen $\sigma$):}
\[
  F(t+\Delta t) = F(t)\exp\!\bigl[-\tfrac{1}{2}\sigma^2\Delta t
                                   +\sigma\sqrt{\Delta t}\,Z_1\bigr].
\]
No discretization error for constant $\sigma$; cross-variation error with
$\sigma$ is $O(\Delta t)$.

\textbf{$\beta=0$ (normal, Euler-Maruyama, exact for constant $\sigma$):}
\[
  F(t+\Delta t) = F(t) + \sigma\sqrt{\Delta t}\,Z_1.
\]

\textbf{$0<\beta<1$ (CEV, Euler with absorbing boundary):}
\[
  F(t+\Delta t) = \max\!\bigl(F(t)+\sigma\max(F(t),0)^\beta\sqrt{\Delta t}\,Z_1,\;0\bigr).
\]

Code: \coderef{FXSABRMCPricer.cs}, method \texttt{StepForward()}.

%%==========================================================================
\section{Hagan SABR Implied Volatility Approximations}
\label{sec:hagan}
%%==========================================================================

\subsection{Lognormal (Black) implied volatility ($\beta$ general)}

Hagan et al.\ (2002) derived an asymptotic expansion for the lognormal (Black)
implied vol $\sigma_B(F,K)$ accurate to $O(T)$.  Define:
\[
  x = \ln(F/K),\quad
  \zeta = \frac{\nu}{\alpha}(FK)^{(1-\beta)/2}x,\quad
  \chi(\zeta) = \ln\!\left(\frac{\sqrt{1-2\rho\zeta+\zeta^2}+\zeta-\rho}{1-\rho}\right).
\]
The implied Black vol is:
\[
  \sigma_B(F,K) = \frac{\alpha}{(FK)^{(1-\beta)/2}\,D(x)}
  \cdot\frac{\zeta}{\chi(\zeta)}\cdot\bigl[1+\mathcal{C}\cdot T\bigr],
\]
where
\[
  D(x) = 1+\frac{(1-\beta)^2}{24}x^2+\frac{(1-\beta)^4}{1920}x^4,
\]
\[
  \mathcal{C} = \frac{(1-\beta)^2\alpha^2}{24(FK)^{1-\beta}}
              + \frac{\rho\beta\nu\alpha}{4(FK)^{(1-\beta)/2}}
              + \frac{(2-3\rho^2)\nu^2}{24}.
\]
Limiting cases:
\begin{itemize}[nosep]
  \item $\nu\to 0$: $\zeta\to 0$, $\zeta/\chi\to 1$; reduces to CEV vol $\alpha/(FK)^{(1-\beta)/2}$.
  \item $F\to K$ (ATM): $x\to 0$, $D\to 1$; no $0/0$ if handled by L'H\^{o}pital.
  \item $\beta=1$, ATM: $\sigma_B\to\alpha(1+\mathcal{C}T)$ with $\mathcal{C}=(\rho\nu\alpha/4+(2-3\rho^2)\nu^2/24)$.
\end{itemize}

Code: \coderef{Hagan.cs}, \texttt{BlackVol()}.

\subsection{Normal implied volatility ($\beta=0$, ATM)}

The Hagan ATM normal-vol formula:
\[
  \sigma_N(F,F) = \alpha\Bigl(1+\frac{(2-3\rho^2)\nu^2}{24}\,T\Bigr).
\]
Combined with the Bachelier formula (Section~\ref{sec:bachelier}), this
benchmarks the normal-SABR ($\beta=0$) MC prices.

Code: \coderef{Hagan.cs}, \texttt{NormalVolAtm()}.

%%==========================================================================
\section{Barrier Options}
\label{sec:barrier}
%%==========================================================================

\subsection{Definitions}

A single-barrier option pays a vanilla payoff $\max(\phi(S(T)-K),0)$ only
if the barrier $H$ is (or is not) crossed by the process during $[0,T]$:
\begin{center}
\begin{tabular}{lll}\toprule
Type & Activated by & If not activated\\\midrule
Up-and-In   & $S(t)\geq H$ for some $t$ & rebate\\
Up-and-Out  & $S(t)< H$ for all $t$ & rebate\\
Down-and-In & $S(t)\leq H$ for some $t$ & rebate\\
Down-and-Out& $S(t)>H$ for all $t$  & rebate\\\bottomrule
\end{tabular}
\end{center}

The model-free identity for zero rebate:
\[
  V_{\rm KI} + V_{\rm KO} = V_{\rm vanilla},
\]
holds path-by-path, providing an exact numerical sanity check (tested in
\coderef{FXBarrierOptionMCPricerTests.cs} and
\coderef{FXSABRBarrierOptionMCPricerTests.cs}).

\subsection{Discrete monitoring}

The barrier is checked at each MC time step $t_k$.  Code:
\coderef{FXBarrierOptionMCPricer.cs}, \texttt{CheckDiscreteBarrier()};
\coderef{FXSABRBarrierOptionMCPricer.cs}, \texttt{CheckDiscreteBarrier()}.

\subsection{Continuous monitoring: Brownian bridge correction}

With discrete time steps the MC pricer may miss barrier crossings occurring
between observation dates.  The Brownian bridge correction (Baldi,
Caramellino, Iovino 1999) estimates the conditional probability that a
continuous GBM path crosses $H$ in $[t_k,t_{k+1}]$, given the endpoints
$a=S(t_k)$ and $b=S(t_{k+1})$ on the \emph{same side} of $H$:

\[
  \boxed{P(\text{cross}\mid a,b,H)
  = \exp\!\left(\frac{-2\ln(H/a)\,\ln(H/b)}{\sigma^2\Delta t}\right).}
\]

\subsubsection{Derivation}

Under the T-forward measure, the log-process $X(t)=\ln F(t)$ is a Brownian
motion with diffusion $\sigma$ (for $\beta=1$).  Conditioned on
$X(t_k)=\ln a$ and $X(t_{k+1})=\ln b$, the bridge $\{X(s):t_k\le s\le t_{k+1}\}$
is a Brownian bridge with variance $\sigma^2(s-t_k)(t_{k+1}-s)/(t_{k+1}-t_k)$.

The maximum of a Brownian bridge $B^{\rm br}$ on $[0,1]$ with $B^{\rm br}(0)=u$,
$B^{\rm br}(1)=v$ satisfies:
\[
  P\!\left(\sup_{s\in[0,1]}B^{\rm br}(s)\geq h\right)
  = \exp\!\bigl(-2h(h-u-v)\bigr)\mathbf{1}_{h>\max(u,v)}.
\]
Translating to the log barrier $h=\ln H$, $u=\ln a$, $v=\ln b$ and
rescaling by total variance $\sigma^2\Delta t$:
\[
  P = \exp\!\left(-2\cdot\frac{(\ln H-\ln a)(\ln H-\ln b)}{\sigma^2\Delta t}\right)
    = \exp\!\left(\frac{-2\ln(H/a)\ln(H/b)}{\sigma^2\Delta t}\right).
\]
When $\ln(H/a)$ and $\ln(H/b)$ have opposite signs (endpoints on opposite
sides), $P=1$ (certain crossing).

Code references:
\begin{itemize}[nosep]
  \item \coderef{FXBarrierOptionMCPricer.cs}: \texttt{BridgeCrossingProb()},
    \texttt{CheckContinuousBarrier()} with $\sigma^2\Delta t=\texttt{SigSqDt}$ (constant vol).
  \item \coderef{FXSABRBarrierOptionMCPricer.cs}: same formula, with
    $\sigma^2\Delta t$ replaced by the stochastic \emph{beginning-of-step} SABR vol:
    \[ \sigma^2\Delta t \;\longrightarrow\; v_{\rm beg}(t)^2\,\Delta t, \]
    where $v_{\rm beg}(0)=\alpha$ and $v_{\rm beg}(t_k)=\texttt{volPath}[k-1]$.
    This is exact for $\beta=1$ (locally lognormal) and an approximation for other $\beta$.
\end{itemize}

\subsubsection{Reproducibility}

Bridge draws use a deterministic seed per (path, step):
\begin{lstlisting}
var pathSeed = CurrentPathIndex ^ unchecked((int)0x9E3779B9u);
var u = new Random(pathSeed ^ (t * 1000003)).NextDouble();
if (u < pCross) return true;
\end{lstlisting}
The salt \texttt{0x9E3779B9} (the fractional part of the golden ratio in
32-bit hex) distributes the seed space.

\subsection{SABR barrier: barrier on the forward}

FX barrier options are quoted on the spot $S$.  The SABR MC pricer simulates
the forward $F(t,T)$ (not $S$), so the barrier is expressed in forward
space.  Under the T-forward measure with flat rates:
\[
  F(t,T) = S(t)\,e^{(r_d-r_f)(T-t)}.
\]
The barrier on $F$ coincides with the barrier on $S$ if and only if $r_d=r_f$
(then $F(t,T)=S(t)$ for all $t$).  This is used in the degeneracy test
\texttt{ZeroVolOfVol\_EqualRates\_SabrBarrier\_MatchesGbmBarrier} in
\coderef{FXSABRBarrierOptionMCPricerTests.cs}.

Code: \coderef{FXSABRBarrierOptionMCPricer.cs}, \texttt{EvaluatePayoff()}.

%%==========================================================================
\section{Monte Carlo Pricing}
\label{sec:mc}
%%==========================================================================

\subsection{Estimator}

Given $M$ independent paths, the Monte Carlo price estimator is:
\[
  \hat V = \frac{1}{M}\sum_{m=1}^M \mathrm{Payoff}_m,
  \qquad \mathrm{Var}[\hat V] = \frac{s^2}{M},
  \quad s^2=\frac{1}{M-1}\sum_m(\mathrm{Payoff}_m-\hat V)^2.
\]
A 95\% confidence interval is $\hat V\pm 1.96\,s/\sqrt{M}$.

Code: \coderef{MCBasePricer.cs}, \texttt{Price()};
\coderef{PricingResult} carries \texttt{Price}, \texttt{StandardError},
\texttt{ConfidenceIntervalLower/Upper}.

\subsection{Parallel path evaluation}

\coderef{MCBasePricer.Price()} uses \texttt{Parallel.For} to distribute
paths across CPU cores with thread-local accumulators (sum, sum-of-squares),
merged at the end.  Thread safety is achieved via:
\begin{itemize}[nosep]
  \item ThreadLocal spot/forward buffers in \coderef{GbmOptionMCPricer.cs}
    and \coderef{FXSABRMCPricer.cs}.
  \item The \texttt{SimulationCube} is read-only after construction.
\end{itemize}

\subsection{Simulation cube: pre-generated standard normals}

The \texttt{SimulationCube} (shape $M\times\text{Steps}\times\text{Assets}$)
pre-generates all required $\N(0,1)$ draws.  For SABR, two assets are
required (one for the forward, one for the independent vol factor).  For
single-asset GBM, one asset suffices.

Code: \coderef{SimulationCube.cs}.

\subsection{Common random numbers (variance reduction)}

Bumped pricers for Greeks reuse the same \texttt{SimulationCube}
(\texttt{CreateBumped} returns a new pricer instance pointing to the same
cube).  The base price and the bumped price share path realizations, reducing
the variance of the finite-difference estimator by positive path-level
correlation.

Code: \coderef{FXSABRVanillaOptionMCPricer.cs} and
\coderef{FXSABRBarrierOptionMCPricer.cs}, method \texttt{CreateBumped()}.

%%==========================================================================
\section{Greeks}
\label{sec:greeks}
%%==========================================================================

\subsection{Bump-and-reprice}

All Greeks are computed by finite-difference bump-and-reprice:
\[
  \Delta = \frac{V(S+\varepsilon)-V(S-\varepsilon)}{2\varepsilon},
  \quad
  \Gamma = \frac{V(S+\varepsilon)-2V(S)+V(S-\varepsilon)}{\varepsilon^2},
  \quad
  \mathcal{V} = \frac{V(\sigma+\varepsilon)-V(\sigma-\varepsilon)}{2\varepsilon}.
\]
For spot/rate bumps the same pricer instance is re-priced with bumped market
data (\coderef{MarketDataBumps.cs}).  For vol bumps, a new pricer instance
is constructed with bumped $\sigma$ or $\alpha$, reusing the same cube.

\subsection{SABR-specific Greeks}

\textbf{Vega ($\partial V/\partial\alpha$):} bump $\alpha\to\alpha\pm\varepsilon$,
construct new pricer, reprice on same market.  $\alpha$ is the SABR initial
vol level, playing the role of ATM vol.

\textbf{Vol-of-vol sensitivity ($\partial V/\partial\nu$):} bump $\nu$,
construct new pricer.  Code: \coderef{FXSABRMCPricer.cs},
\texttt{ComputeVolOfVolSensitivity()}.

\textbf{Correlation sensitivity ($\partial V/\partial\rho$):} bump $\rho$
(clamped to $(-1,1)$), construct new pricer.  Code:
\coderef{FXSABRMCPricer.cs}, \texttt{ComputeCorrelSensitivity()}.

\textbf{Theta:} not computed (returned as \texttt{NaN}) because bumping time
would require rebuilding the cube or re-parameterising the contract.

%%==========================================================================
\section{Basket Options}
\label{sec:basket}
%%==========================================================================

\subsection{Model}

An $n$-leg basket with weights $w_i$ and a weighted-sum aggregation:
\[
  B(t) = \sum_{i=1}^n w_i\,S_i(t).
\]
Each leg $S_i$ follows correlated GBM under $\Q$:
\[
  \frac{\dd S_i}{S_i} = \mu_i\,\dd t + \sigma_i\,\dd W_i,
  \qquad \dd W_i\,\dd W_j = \rho_{ij}\,\dd t.
\]
The simulation uses a Cholesky decomposition of the correlation matrix to
generate correlated standard normals from independent draws.

Code: \coderef{FXBasketVanillaMCPricer.cs},
\coderef{FXBasketBarrierOptionMCPricer.cs}.

\subsection{Effective vol for basket barrier bridge}

For continuous barrier monitoring on a basket, the effective basket
volatility $\sigma_{\rm eff}$ at each step is reconstructed from the realized
leg spots and the leg vols, used as the local diffusion coefficient in the
Brownian bridge formula (Section~\ref{sec:barrier}).

%%==========================================================================
\section{Normal Distribution Utilities}
\label{sec:normal}
%%==========================================================================

The standard normal CDF is evaluated using the complementary error function:
\[
  \Phi(x) = \tfrac{1}{2}\,\mathrm{erfc}\!\bigl(-x/\sqrt{2}\bigr),
\]
via \texttt{MathNet.Numerics.SpecialFunctions.Erfc}, giving full
double-precision accuracy (no polynomial approximation error).

The PDF:
\[
  \phi(x) = \frac{1}{\sqrt{2\pi}}\,e^{-x^2/2}.
\]

Code: \coderef{NormalDistribution.cs}, \texttt{Cdf()}, \texttt{Pdf()}.

%%==========================================================================
\section{Summary of Code-to-Formula Mapping}
\label{sec:summary}
%%==========================================================================

\begin{center}
\renewcommand{\arraystretch}{1.25}
\begin{tabular}{lll}\toprule
\textbf{File} & \textbf{Model / Formula} & \textbf{Section}\\\midrule
\texttt{GarmanKohlhagen.cs}           & GK call/put                       & \ref{sec:gk}\\
\texttt{BlackScholes.cs}              & BSM call/put                      & \ref{sec:bsm}\\
\texttt{BlackForward.cs}              & Black forward call/put            & \ref{sec:black_fwd}\\
\texttt{Bachelier.cs}                 & Normal model call/put             & \ref{sec:bachelier}\\
\texttt{DigitalBlackScholes.cs}       & Cash/asset-or-nothing             & \ref{sec:digital}\\
\texttt{GeometricAsian.cs}            & Kemna-Vorst geometric Asian       & \ref{sec:asian}\\
\texttt{Hagan.cs}                     & SABR $\sigma_B$, $\sigma_N$       & \ref{sec:hagan}\\
\texttt{NormalDistribution.cs}        & $\Phi$, $\phi$ via erfc           & \ref{sec:normal}\\
\texttt{ZeroCurve.cs}                 & Interpolation, $P(0,T)$, $f(t_0,t_1)$ & \ref{sec:curves}\\
\texttt{GbmOptionMCPricer.cs}         & GBM log-Euler, discrete dividends & \ref{sec:gbm},\ref{sec:dividends}\\
\texttt{FXSABRMCPricer.cs}            & SABR under $\QT$, Euler-Maruyama  & \ref{sec:sabr}\\
\texttt{FXBarrierOptionMCPricer.cs}   & Barrier, Brownian bridge (GBM)    & \ref{sec:barrier}\\
\texttt{FXSABRBarrierOptionMCPricer.cs}&Barrier, Brownian bridge (SABR)   & \ref{sec:barrier}\\
\texttt{FXBasketVanillaMCPricer.cs}   & Correlated GBM basket             & \ref{sec:basket}\\
\texttt{FXBasketBarrierOptionMCPricer.cs}&Basket barrier + bridge         & \ref{sec:basket},\ref{sec:barrier}\\
\texttt{MCBasePricer.cs}              & MC estimator, confidence interval & \ref{sec:mc}\\\bottomrule
\end{tabular}
\end{center}

%%==========================================================================
\section*{References}
%%==========================================================================
\addcontentsline{toc}{section}{References}

\begin{enumerate}[label={[\arabic*]}]
  \item Garman, M.\ and Kohlhagen, S.\ (1983).  ``Foreign currency option values.''
        \textit{Journal of International Money and Finance} 2(3): 231--237.
  \item Black, F.\ and Scholes, M.\ (1973).  ``The pricing of options and corporate liabilities.''
        \textit{Journal of Political Economy} 81(3): 637--654.
  \item Black, F.\ (1976).  ``The pricing of commodity contracts.''
        \textit{Journal of Financial Economics} 3(1--2): 167--179.
  \item Bachelier, L.\ (1900).  ``Th\'eorie de la sp\'eculation.''
        \textit{Annales Scientifiques de l'\'Ecole Normale Sup\'erieure} 17: 21--86.
  \item Hagan, P., Kumar, D., Lesniewski, A., Woodward, D.\ (2002).
        ``Managing smile risk.''
        \textit{Wilmott Magazine}: 84--108.
  \item Baldi, P., Caramellino, L., Iovino, M.G.\ (1999).
        ``Pricing general barrier options: A numerical approach using sharp large deviations.''
        \textit{Mathematical Finance} 9(4): 293--321.
  \item Kemna, A.\ and Vorst, A.\ (1990).
        ``A pricing method for options based on average asset values.''
        \textit{Journal of Banking and Finance} 14(1): 113--129.
  \item Merton, R.C.\ (1973).  ``Theory of rational option pricing.''
        \textit{Bell Journal of Economics and Management Science} 4(1): 141--183.
  \item Brigo, D.\ and Mercurio, F.\ (2006).
        \textit{Interest Rate Models -- Theory and Practice}, 2nd ed.\ Springer.
\end{enumerate}

\end{document}
